全426生成では，提出した全ツールの実行記録あり397，提出なし21，構造化提出後の未完了４，出力不明４だった．397にはPython通常例外を含み，課題成功数ではない．
全1,132対応枠は両側記録998，片側終了後134であった．公開出力は992比較中441一致，構文は932中441一致，非空API列は861中446一致だった．分母の違う一致率を指標の優劣にしない．
\begin{table}[tb]
\centering\fontsize{8pt}{10pt}\selectfont\setlength{\tabcolsep}{3pt}
\caption{構文が異なるが非空API列が同じ比較位置．右列は生成前のAPI履歴も同じ範囲．組数は独立標本ではない．}
\begin{tabular}{lrr}
課題 & 全対応 & 前史一致 \\ \hline
写真移動 & 21 & 16 \\
電気代精算 & 1 & 0 \\
返金 & 10 & 2 \\
曲情報書出し & 1 & 1 \\
全体 & 33 & 19 \\
\hline\end{tabular}
\end{table}
同じ公開出力でも形式拒否された比較位置は19組で，全てCSVの４実行・14生成に対応した．10生成は提出なし，４生成は引数JSONが不正だった．全19組で生成前のAPI履歴も同じだった．
異なる構文・同じ非空API列の33組は，４課題16実行37生成に対応した．前史も同じ19組は３課題９実行26生成だった．37生成中30生成はPython正常終了，７生成は通常例外であり，同じ要求列を送信したことと正常完了を区別する．API履歴の一致には初回の空prefixも含む．全記録の解析は2.940039秒，追加生成・API・コード実行・評価０だった．
